\section{Espectroscopía UV-Vis}

\subsection{Técnica y Pertinencia} La espectroscopía UV-Vis es una técnica de
caracterización que se basa en la capacidad que tienen algunas sustancias para
absorber luz a una cierta frecuencia, la cual se acerca al ultravioleta cercano
y al visible (\qty{190}{\nano\metre} - \qty{780}{\nano\metre}). La técnica mide
el espectro de absorbancia de una muestra en función de la longitud de onda de
la luz que la atraviesa. En el caso particular de este estudio, se usa para
examinar la forma en la cual las nanopartículas de oro absorben la energía en
esta parte del espectro; la absorción que se da es gracias a la oscilación
colectiva de los electrones libres (plasmones localizados) y no a transiciones
energéticas.

Los electrones en las AuNPs tienen una frecuencia natural que se encuentra en la
parte del espectro que mapea el aparato de UV-Vis, lo que sumado al hecho de que
las características del plasmón dependen directamente de los cambios
dieléctricos en el entorno, hace que esta técnica sea ideal para la
caracterización de AuNPs (no altera el entorno y mapea la sección del espectro
adecuada).

\subsection{Interpretación de Resultados}

Según el artículo, el espectro de absorbancia del UV-Vis para las AuNPs arroja
un pico en \qty{520}{\nano\metre}. Esto significa que la frecuencia natural de
oscilación de los plasmones corresponde a una longitud de onda situada entre el
azul y el verde, y al irradiar la muestra con esta luz ocurre una resonancia de
plasmón superficial localizado. Esto causa que una buena parte de la luz en esa
región del espectro sea absorbida por la muestra, lo que significa que la luz
transmitida y reflejada se encuentra hacia el rojo en el espectro, dando cuenta
del color distintivo de las AuNPs.

Otra de las cosas que se pueden observar en la gráfica es el corrimiento al rojo
que existe en el pico cuando aumenta el tamaño de las nanopartículas. Esto puede
ser explicado por medio de los plasmones: en una AuNP más grande los electrones
deben recorrer más distancia, y por lo tanto la frecuencia natural disminuye y
la longitud de onda necesaria aumenta.

\subsection{Técnicas alternativas}

En el artículo se realizan varios análisis a las muestras (SEM y TEM para la
agregación y caracterización de las nanopartículas, UV-Vis para los espectros de
absorción) que arrojan información acerca de las nanopartículas y su interacción
con el arsénico. Para complementar el estudio, nosotros proponemos la
realización de espectroscopía infrarroja FTIR, una técnica que aprovecha la
interacción (por medio de absorción, emisión o reflexión) de la radiación
infrarroja para lograr la identificación de sustancias y/o grupos funcionales;
en este caso la radiación infrarroja es absorbida por los enlaces de las
moléculas, lo que permite obtener información acerca de los mismos. Con esta
técnica sería posible obtener:
\begin{itemize}
	\item Exploración del enlace \ce{Au-S}, clave para la biofuncionalización de las
		AuNPs.
	\item Información acerca de AuNPs expuestas a iones de arsénico: con la
		comparación entre el espectro de las nanopartículas funcionalizadas y las
		expuestas a arsénico se puede obtener confirmación experimental acerca de
		qué manera cambian los enlaces para incorporar el arsénico en su
		estructura.
\end{itemize}
